\subsection{集合论}

集合论是一种包含关系符号的理论 \(=,\) \(\in\) 以及实质性的符号 \(\supset\) (所有这些符号的权均为2)；除了第一章中给出的模式S1至S7外，它还包含将在第6节中引入的模式S8，以及显式公理A1（第3节）、A2（第5节）、A3（ \(\S\) 2，第1号），A4（ \(\S\) 5，第1期），以及A5（第三章， \(\S\) 6，第1期），这些显式公理不包含任何字母；换言之，集合论是一个没有常量的理论。

由于集合论是一种平等理论，第一章的结果适用。

从现在起，除非另有明确说明，我们总是在比集合论更强的理论（第一章，§2，第4节）中进行论证；如果未提及该理论，则假定隐含集合论。在许多情况下，这一假设显然是多余的，读者应不难确定所陈述的结果在何种弱于集合论的理论中成立。

如果T和U是项，则这个汇编 \(\in TU\) 是一个关系（称为属于关系），在实际中我们通常以下列方式之一书写： \(\widetilde{T} \in U\) , \((\widetilde{T}) \in (U)\) “T 属于 U”，“T 是 U 的一个元素”。关系“不 \((\widetilde{T} \in U)\) '' 写作 \(T \notin U\) .

从“朴素”的观点来看，许多数学实体可以被视为对象的集合或“集”。我们不寻求将这一概念形式化，在后续内容的形式主义解释中，“集”一词应严格视为“项”的同义词。特别是，诸如“设X为一个集”这样的短语，原则上完全是多余的，因为每个字母都是一个项。引入这些短语仅是为了帮助对文本的直观理解。

